\section*{ВВЕДЕНИЕ}
\addcontentsline{toc}{section}{ВВЕДЕНИЕ}

Современные города активно развивают цифровые сервисы для жителей и туристов. Туристическая отрасль нуждается в удобных инструментах, которые позволяют быстро ориентироваться в городском пространстве: находить достопримечательности, строить маршруты, получать актуальную информацию о местах и событиях. При отсутствии специализированной системы сведения о городских объектах часто разрознены: описания хранятся в текстовых документах, фотографии — в социальных сетях, а маршруты приходится прокладывать с помощью внешних картографических сервисов. Такой подход приводит к дублированию данных, сложности их актуализации и снижению доверия пользователей.

Сайт-гид является лицом города в цифровом пространстве и может существенно повысить его привлекательность для туристов. Любой пользователь сети Интернет сможет получить необходимую информацию о достопримечательностях, спланировать прогулку и найти контактные данные организаций в любой момент. Поэтому создание удобного, современного и информативного туристического портала становится важной задачей.

\emph{Цель настоящей работы} --- разработка веб-приложения «Цифровая платформа Курский край», предоставляющего пользователям интерактивный гид по городу Курску с возможностью просмотра мест, фильтрации по категориям, построения пешеходных маршрутов, фотогалереи и административной панели для управления контентом. Для достижения поставленной цели необходимо решить \emph{следующие задачи:}
\begin{itemize}
\item провести анализ предметной области и существующих решений в сфере городских туристических сервисов;
\item разработать концептуальную модель веб-сайта и определить требования к системе;
\item спроектировать архитектуру клиентской и серверной частей, базу данных и пользовательский интерфейс;
\item реализовать веб-приложение с использованием современных веб-технологий.
\end{itemize}

\emph{Структура и объем работы.} Отчет состоит из введения, 4 разделов основной части, заключения, списка использованных источников, 2 приложений. Текст выпускной квалификационной работы равен \formbytotal{lastpage}{страниц}{е}{ам}{ам}.

\emph{Во введении} сформулирована цель работы, поставлены задачи разработки, описана структура работы, приведено краткое содержание каждого из разделов.

\emph{В первом разделе} на стадии описания технической характеристики предметной области приводится сбор информации о деятельности компании, для которой осуществляется разработка сайта.

\emph{Во втором разделе} на стадии технического задания приводятся требования к разрабатываемому сайту.

\emph{В третьем разделе} на стадии технического проектирования представлены проектные решения для web-сайта.

\emph{В четвертом разделе} приводится список классов и их методов, использованных при разработке сайта, производится тестирование разработанного сайта.

В заключении излагаются основные результаты работы, полученные в ходе разработки.

В приложении А представлен графический материал.
В приложении Б представлены фрагменты исходного кода. 
